\chapter*{Resumen}
\addcontentsline{toc}{chapter}{Resumen}

Los centros de operaciones de seguridad (SOC) se enfrentan hoy a dos problemas que se
refuerzan mutuamente. Por un lado, el volumen de alertas generadas por las herramientas
de detección supera con holgura la capacidad de análisis humano, lo que provoca el
fenómeno conocido como \textit{fatiga de alertas}: los analistas descartan o posponen
avisos que podrían ser relevantes simplemente porque no hay tiempo material para
revisarlos todos. Por otro, esas alertas llegan \textit{sin contexto}: un aviso aislado no
indica si el comportamiento observado corresponde a una técnica documentada, si ha sido
empleada por un actor conocido, ni qué activos de la organización quedan expuestos como
consecuencia.

A este panorama se suma que los ataques dirigidos no se manifiestan como eventos puntuales,
sino como \textbf{secuencias de etapas} que se desarrollan a lo largo del tiempo:
reconocimiento, acceso inicial, ejecución, persistencia, movimiento lateral y exfiltración.
Los sistemas de detección que evalúan cada evento de forma independiente son, por
construcción, incapaces de reconocer esa estructura temporal, y resultan por tanto fáciles
de evadir mediante técnicas individualmente poco llamativas.

Este Trabajo Fin de Grado propone el diseño, la implementación y la validación de una
plataforma de datos que aborda ambos problemas de forma conjunta. El sistema ingiere
telemetría heterogénea de red y de sistema operativo, la normaliza a un esquema canónico
(OCSF), detecta secuencias de ataque multi-etapa mediante procesamiento de eventos
complejos (CEP), las contextualiza contra un grafo de conocimiento de inteligencia de
amenazas alineado con MITRE ATT\&CK, razona sobre las rutas de ataque factibles hacia los
activos críticos empleando un motor lógico basado en Datalog, y emite finalmente una
priorización de incidentes \textbf{explicable}, en la que cada puntuación va acompañada de
la descomposición de los factores que la justifican.

La arquitectura se articula sobre tres pilares: la detección multi-etapa, el grafo de
ataque razonado y la priorización explicable. El objetivo último no es sustituir a las
herramientas de detección existentes, sino transformar su salida en información accionable,
aportando \textbf{conciencia cibersituacional} y soporte a la decisión al analista.

La validación se realiza sobre conjuntos de datos públicos etiquetados con MITRE ATT\&CK,
lo que permite contrastar cuantitativamente la precisión de la detección, la reducción del
volumen de alertas y la calidad de la priorización frente a líneas base establecidas.

\vspace{1em}
\noindent\textbf{Palabras clave:} conciencia cibersituacional, detección de intrusiones,
MITRE ATT\&CK, procesamiento de eventos complejos, grafos de ataque, inteligencia de
amenazas, ingeniería de datos en tiempo real, explicabilidad.
